\section{What it cost, and what it caught}

The previous section argued that the signatures of a working method arrive disguised as friction. This one puts numbers to the friction, and then makes a claim that follows from \S4.8 and is easy to miss: \textbf{the costs and the catches are not two lists. They are the same list, read twice.}

\subsection{What the friction cost}

\textbf{Verification costs an iteration.} Checking a document's claims at source before typesetting it took a full cycle and produced no prose. Re-scoring every recorded verdict under a corrected criterion took another. Two of thirteen experiments' worth of effort produced no new science.

\textbf{Pre-registration costs a commit and an argument.} Declaring what would falsify a hypothesis, and separately what would mean the experiment asked the wrong question, before the runner exists --- that is genuine work done at the moment when the answer is least visible and the temptation to defer is strongest.

\textbf{Rejections cost round trips.} Three specifications were rejected on false premises. Each cost a cycle.

\textbf{The record costs reading.} Nine thousand lines of ledger in five working days, and one human reading every pull request.

\subsection{What the friction caught}

The same events.

The verification iteration found three wrong figures and eight claims true only in the condition they were measured under. The re-scoring iteration found that a criterion had governed ten experiments without testing what it was read as testing. The three rejections each stopped work that would have measured the wrong thing --- one of them a premise so central that the experiment would have answered a different question entirely and reported the answer as if it were the one asked.

\textbf{A wrong belief costs everything downstream of it.} That is not rhetoric; it is the arithmetic of this project. The criterion defect propagated through ten experiments and required re-reading every null result they had produced. Caught at experiment two, it would have cost an afternoon.

So the honest accounting is not \textit{the method cost two iterations and caught four errors.} It is: \textbf{the method spent two iterations to avoid re-running ten.}

\subsection{The cost of success}

The costs above are the ones a reader expects. These are the ones that surprised, and every one of them arrives when things are going \textit{well}.

\textbf{A positive result creates pressure to build on it.} The project's one positive finding was a single experiment old, underpowered in its most important band, and untested against its own dose-response. The obvious next move was to build the next stage on it. The consolidation done instead found the criterion defect --- which is to say, \textbf{the discipline that paid best was the one applied at the moment of greatest momentum.} Had the result been negative, nobody would have wanted to build on it and the consolidation would have happened for free.

\textbf{A run of clean iterations is when framing stops being examined.} \S4.7's trigger. Worth restating as a cost rather than a mechanism: the better a project is going, the more expensive its next unexamined assumption becomes, because more has been built on top before anyone checks.

\textbf{Rigour is self-reinforcing, and that is not entirely good.} Every artifact that meets a high standard raises the felt cost of producing one that does not. The low-ceremony lane --- where an idea can be tried and thrown away --- shipped on day one and was never used once. \S6 argues this is the method's deepest problem and not an oversight.

\textbf{The record outgrows the reader.} A methodology that produces more record than anyone can read has moved the bottleneck rather than removed it. In the predecessor project, two failures got in exactly there, under merge pressure, with green tests. The apparatus catches wrong \textit{numbers} well and wrong \textit{merges} poorly, and the variable that moved was the rate.

\textbf{The deliverable is unfashionable.} The better this works, the more of the output takes the form of \textit{a well-ordered set of eliminated possibilities} --- harder to publish than a benchmark number, and harder to describe at a conference. That is a fact about publishing rather than about the knowledge, but it is a real cost borne by whoever adopts the method.

\subsection{The asymmetry that justifies the whole thing}

One number carries the argument. The \textbf{first} of the two verification passes over the project's technical report checked twenty-two claims, eighteen of them substantive and four bibliographic. Of those eighteen, \textbf{seven stood as written.}

The other eleven were not fabrications. They were claims made by someone who had read the record, believed them, and would have published them. Three were wrong on the figure; eight were true in the condition they came from and misleading without it.

That is the ratio the method exists for. In software, shipping eleven of eighteen statements slightly wrong is a bad release. In research, it is a body of work that other people build on.

The costs are all paid in time. The failure is paid in credibility, and it is paid by everyone who cited you.
